% PURPOSE: Introduction (budget 1.0 page).
% WHAT GOES HERE: Human-written prose only, replacing the paragraph-job comments below.
% FILLED IN AT STEP: W5.2; evening 2 (Thursday 22 October), reread on evening 8.
% DO NOT: Paste AI-generated text. Do not cite anything without a row in references/citations.md.
\section{Introduction}
% P1 job: VTC works: text sent to a model as an image, to spend fewer tokens. Evidence: cite Glyph and DeepSeek-OCR (rows in citations.md).
% P2 job: the open question: which rendering choices change what a VLM transcribes, and by how much; rendering is inherited from design for human eyes. Evidence: cite LensVLM; FreeType hinting docs (needs a row in citations.md).
% P3 job: what prior work shows at coarse scale, then the gap: glyph-level, preregistered measurement across readers. Evidence: cite Grid2Matrix, Steal the Patch Size, CodeOCR.
% P4 job: what we did: three studies, five readers, 40,788 transcriptions, two preregistrations, a frozen pipeline [T]. Evidence: row counts per study in LAB_NOTEBOOK.md; the tags in HANDOFF.md.
% P5 job: the map in brief (headline C): the order of the levers; the largest and the smallest, with their sizes; which rows are preregistered and which exploratory. Evidence: the levers table (section 7).
% P6 job: contributions as a list, each mapped to a section: the map of levers (sections 5 to 7); the spaces mechanism and its preregistered test (sections 5 and 6); a calibration that prices text in vision tokens, with its page-load caveat (section 7); the written specification and the frozen pipeline (section 4).
% P7 job (or the end of P6): your one sharp takeaway for builders, the same as the abstract's.
